\documentclass[../main]{subfiles}
\begin{document}
\newpage

\subsubsection{\texttt{armax} --- ARMAX}

\begin{lstlisting}
theta, m = pysib.armax(u, y, na, nb, nc, nz)
\end{lstlisting}

Identifies an ARMAX model by minimizing the prediction-error
criterion.  The model structure is

\[
  A(q^{-1})\,y(t) =
  B(q^{-1})\,u(t - n_k) + C(q^{-1})\,e(t),
\]

where $e(t)$ is white noise.
The one-step-ahead prediction error is

\[
  \varepsilon(t;\theta) =
  \frac{A(q^{-1})}{C(q^{-1})}\,y(t)
  - \frac{B(q^{-1})}{C(q^{-1})}\,u(t-n_k),
\]

and the criterion $V(\theta) = \tfrac{1}{N}\sum_t \varepsilon(t;\theta)^2$
is minimized.  Because $C$ appears in the denominator, the predictor
depends on past residuals and the problem is nonlinear in $\theta$.
An ARX estimate (with $C$ initialized to zero) provides the initial
guess; the optimizer uses steepest-descent and Newton steps
implemented in a C extension, with LAPACK used for the Newton step.

\subsubsection*{Arguments}
\begin{center}
\begin{tabular}{lll}
  \toprule
  Name & Type & Description \\
  \midrule
  \texttt{u}  & array\_like & Input signal. \\
  \texttt{y}  & array\_like & Output signal. \\
  \texttt{na} & int & Number of $A$ parameters. \\
  \texttt{nb} & int & Number of $B$ parameters. \\
  \texttt{nc} & int & Number of $C$ parameters. \\
  \texttt{nz} & int & Input delay $n_k \geq 0$. \\
  \bottomrule
\end{tabular}
\end{center}

\subsubsection*{Returns}
\begin{center}
\begin{tabular}{lll}
  \toprule
  Name & Type & Value \\
  \midrule
  \texttt{theta} & ndarray & $[a_1,\ldots,a_{n_a},\; b_0,\ldots,b_{n_b-1},\; c_1,\ldots,c_{n_c}]$ \\
  \texttt{m}     & dict    & \texttt{A}, \texttt{B}, \texttt{C} free; \texttt{D}$=$\texttt{F}$=[1]$ \\
  \bottomrule
\end{tabular}
\end{center}

\end{document}
